% Pollis drill-down: Compose > Orchestration, coding agents and the models they run on (Model Tasks and Routing, local
% models through llama.cpp, what happens when a task starts), and the pages where agents work together on research
% (14 November 2026). For users who want the detail, and for whoever maintains Pollis.
% Compile with: latexmk -pdf workflow_orchestration.tex
% Shared look for the one-page Pollis workflow handouts (same fonts and colours as the brochure).
\documentclass[10pt,a4paper]{article}

\usepackage[a4paper,left=17mm,right=17mm,top=14mm,bottom=16mm,footskip=8mm]{geometry}
\usepackage[T1]{fontenc}
\usepackage[default,semibold]{sourcesanspro}
\usepackage{microtype}
\usepackage{graphicx}
\usepackage[table]{xcolor}
\usepackage{tikz}
\usepackage[most]{tcolorbox}
\usepackage{fontawesome5}
\usepackage{enumitem}
\usepackage{fancyhdr}
\usepackage[hidelinks]{hyperref}

\graphicspath{{images/}}

% ---------------------------------------------------------------- colours (as in the brochure)
\definecolor{navy}{HTML}{1C3A5E}
\definecolor{pblue}{HTML}{3B78B5}
\definecolor{amber}{HTML}{F2B33D}
\definecolor{deepamber}{HTML}{D9902A}
\definecolor{pale}{HTML}{EEF3F9}
\definecolor{ink}{HTML}{2A2F36}
\definecolor{logonavy}{HTML}{173D63}
\definecolor{logodark}{HTML}{C98A1E}
\definecolor{logoamber}{HTML}{F2B446}
\definecolor{logocream}{HTML}{FEEFD2}
\definecolor{soft}{HTML}{5E6773}
\color{ink}

\setlength{\parindent}{0pt}
\setlength{\parskip}{3pt}
\setlist{nosep,leftmargin=*,itemsep=2pt}

\pagestyle{fancy}
\fancyhf{}
\renewcommand{\headrulewidth}{0pt}
\fancyfoot[L]{\footnotesize\color{soft}\textbf{Pollis}\enspace Bright Insights through Exploration, Modelling, Simulation and Optimisation}
\fancyfoot[R]{\footnotesize\color{soft}github.com/saragga/pollis/releases/latest}

% The Pollis logo, redrawn so it prints sharply at any size
\newcommand{\pollislogo}[1]{%
  \begin{tikzpicture}[x=#1/512,y=-#1/512,baseline=(current bounding box.center)]
    \fill[logonavy] (0,0) rectangle (512,512);
    \fill[logodark] (36,422) -- (104,256) -- (172,422) -- cycle;
    \fill[logoamber] (104,422) -- (203,180) -- (302,422) -- cycle;
    \fill[logocream] (203,422) -- (339,89) -- (476,422) -- cycle;
  \end{tikzpicture}}

% Top of the page: logo, name, kind of document (Workflow unless given, e.g. [Drill-Down]) and its subject;
% then the question and a one-line promise
\newcommand{\workflowhead}[4][Workflow]{%
  \thispagestyle{fancy}%
  \begin{minipage}[c]{13mm}\pollislogo{11mm}\end{minipage}%
  \begin{minipage}[c]{\dimexpr\linewidth-13mm\relax}
    {\color{navy}\Large\bfseries Pollis}\enspace{\color{soft}#1}\hfill
    {\color{deepamber}\small\bfseries\MakeUppercase{#2}}
  \end{minipage}\par\vspace{7mm}
  {\color{navy}\fontsize{22}{26}\selectfont\bfseries #3}\par\vspace{4pt}
  {\color{soft}\large #4}\par\vspace{5mm}}

% A menu path: \menu{Explore}{Access Data Libraries}... is written as \path{Explore \gt Access Data Libraries}
\newcommand{\gt}{\,{\color{deepamber}\bfseries ›}\,}
\newcommand{\menupath}[1]{{\footnotesize\bfseries\color{navy}#1}}

% A screenshot if images/<file> exists, otherwise an orange placeholder that says what to capture
\newlength{\shotheight}\setlength{\shotheight}{43mm}
\newcommand{\stepshot}[2]{%
  \IfFileExists{images/#1}{%
    \tcbox[enhanced, colback=black, colframe=black!70, boxrule=0.4pt, arc=2pt,
      left=0pt,right=0pt,top=0pt,bottom=0pt, boxsep=0pt, halign=center]{%
      \includegraphics[width=\dimexpr\linewidth-1pt\relax,height=\shotheight,keepaspectratio]{#1}}}{%
    \begin{tcolorbox}[enhanced, colback=white, colframe=orange!80, boxrule=0.6pt, arc=2pt,
      height=\shotheight, valign=center, halign=center, left=6pt, right=6pt]
      {\color{orange}\Large\faCamera}\par\smallskip
      {\color{orange}\footnotesize\bfseries Screenshot: #2}\par\smallskip
      {\color{orange!80!black}\scriptsize images/#1}
    \end{tcolorbox}}}

% One step: number, title, menu path, image file, what to capture, what it means
\newcommand{\step}[6]{%
  \begin{tcolorbox}
    {\color{deepamber}\Large\bfseries #1}\enspace{\color{navy}\large\bfseries #2}\par\vspace{1pt}
    \menupath{#3}\par\vspace{4pt}
    \stepshot{#4}{#5}\par\vspace{3pt}
    {\small #6}
  \end{tcolorbox}}

\newenvironment{steps}{%
  \begin{tcbraster}[raster columns=2, raster equal height=rows, raster column skip=4mm,
    raster row skip=4mm, enhanced, colback=pale, colframe=pale, boxrule=0pt, sharp corners,
    left=6pt,right=6pt,top=6pt,bottom=6pt, fontupper=\raggedright]}{\end{tcbraster}}

% Boxes for the bottom of the page
\newtcolorbox{pbox}[2][]{enhanced, colback=pale, colframe=amber, boxrule=0pt,
  leftrule=3pt, sharp corners, left=7pt, right=7pt, top=6pt, bottom=6pt,
  fontupper=\small\raggedright, before upper={{\color{navy}\normalsize\bfseries #2}\par\smallskip}, #1}
\newtcolorbox{nbox}[1][]{enhanced, colback=navy, colframe=navy, boxrule=0pt,
  sharp corners, left=8pt, right=8pt, top=6pt, bottom=6pt, coltext=white,
  fontupper=\small\raggedright, #1}

% A Next Steps item inside the bottom boxes
\newcommand{\nsitem}[2]{\item[\color{deepamber}\faChevronRight] \textbf{#1}: #2}

% The page of each heading, for Ask Pollis
% Pollis: the page of each heading, for Ask Pollis. A document that inputs this file, at the end of its preamble, writes
% <jobname>.labels beside its PDF, a line page|level|heading for each heading, the page counted from the first, the
% cover, as 1. compile_documents.jl copies it to ~/.pollis/documents, and ask_index.jl cuts the document into passages
% under its headings with it. \asklabel{1}{Title} marks a topic or a page, \asklabel{2}{Title} a heading inside it; the
% sections and subsections of a document that has them are marked as they go into its contents. A package vignette
% inputs it before its \documentclass (compile_vignette.jl), so the contents are taken over when the document begins,
% after the packages have taken them over.
\makeatletter
\newwrite\pollis@labels
\immediate\openout\pollis@labels=\jobname.labels
\newcommand{\asklabel}[2]{\write\pollis@labels{\the\ReadonlyShipoutCounter|#1|\unexpanded{#2}}}
\def\pollis@toc{toc}
\newcommand{\pollis@contentsline}[3]{%
  \def\pollis@file{#1}%
  \ifx\pollis@file\pollis@toc\asklabel{#2}{#3}\fi
  \pollis@addcontentsline{#1}{#2}{#3}}
% The headings of the New Methods style of the vignettes, \nmsection and \nmsub, which go into no contents
\def\pollis@asksection#1{\pollis@nmsection{#1}\asklabel{1}{#1}}
\def\pollis@asksub#1{\pollis@nmsub{#1}\asklabel{2}{#1}}
\AddToHook{begindocument/end}{%
  \let\pollis@addcontentsline\addcontentsline\let\addcontentsline\pollis@contentsline
  \ifdefined\nmsection\let\pollis@nmsection\nmsection\let\nmsection\pollis@asksection\fi
  \ifdefined\nmsub\let\pollis@nmsub\nmsub\let\nmsub\pollis@asksub\fi}
\makeatother



\usepackage{titlesec}
\usepackage{tabularx}
\usepackage{booktabs}
\usepackage{listings}

\fancyhf{}
\fancyfoot[L]{\footnotesize\color{soft}\textbf{Pollis}\enspace Drill-down: Orchestration}
\fancyfoot[R]{\footnotesize\color{soft}\thepage}

\titleformat{\section}{\color{navy}\large\bfseries}{\color{deepamber}\thesection}{0.6em}{}
\titleformat{\subsection}{\color{navy}\normalsize\bfseries}{\color{deepamber}\thesubsection}{0.5em}{}
\titlespacing*{\section}{0pt}{12pt}{4pt}
\titlespacing*{\subsection}{0pt}{8pt}{2pt}
\setlength{\parskip}{4pt}

\lstset{basicstyle=\ttfamily\footnotesize, columns=fullflexible, keepspaces=true, breaklines=true,
  backgroundcolor=\color{pale}, frame=leftline, rulecolor=\color{amber}, framerule=2pt, xleftmargin=6pt,
  framexleftmargin=6pt, aboveskip=4pt, belowskip=4pt, upquote=true, tabsize=2}

\usepackage{xurl}
\newcommand{\file}[1]{{\small\nolinkurl{#1}}}
\newcolumntype{L}{>{\raggedright\arraybackslash}X}
\newcommand{\ui}[1]{\textbf{#1}}
\newcommand{\setting}[1]{\file{#1}}
% Code with spaces in it, such as a command line, which \file would run together
\newcommand{\code}[1]{\texttt{\small #1}}
% A menu path at the start of a table cell
\newcommand{\cellpath}[1]{\leavevmode\menupath{#1}}

\begin{document}

\workflowhead[Drill-Down]{Orchestration}
  {Coding agents and their models}
  {Which agent and which model run each task of the Compose menu, how a local model takes the place of a cloud one,
   what happens when a task starts, and the pages where several agents work on one piece of research.}

\begin{nbox}
{\color{amber}\bfseries Drill-down.}\enspace
A look under the surface of what the Guide describes in \emph{Orchestration}. It is for anyone who wants to choose
the models of their tasks with care, run them on models of their own computer, or know what the research pages of
\textbf{Compose \gt Orchestration} will do; \S\ref{sec:machinery} is for whoever maintains Pollis. It describes the
state of 14 November 2026.
\end{nbox}

\tableofcontents

\section{What Orchestration is}\label{sec:what}
A \emph{coding agent} is a program that writes code and documents in a terminal, on a language model of a provider's
cloud or of your own computer. Pollis gives agents the tasks of the \menupath{Compose} menu that need one: creating
panels, packaging toolboxes, writing methods and vignettes, replicating articles, polishing reports. Orchestration
decides which agent, on which model, does each task, and lets several agents work on one piece of research.

\begin{center}\small
\renewcommand{\arraystretch}{1.25}
\begin{tabularx}{\linewidth}{@{}p{0.36\linewidth}L@{}}
\toprule
\textbf{Menu item} & \textbf{What it is for}\\
\midrule
\cellpath{Compose \gt Coding Agents} & The agents: Claude Code, Codex, OpenCode and your own, whether each is on this
  computer, and the Pollis skills installed in them\\
\cellpath{Compose \gt Agent Skills} & The skills: the know-how an agent reads to do a task\\
\cellpath{Compose \gt Orchestration \gt Model Tasks and Routing} & Which agent, on which model, runs each task\\
\cellpath{Compose \gt Local AI Models} & The models of your own computer, which a task can run on\\
\cellpath{Compose \gt Orchestration}, the research pages & Several agents on one question: Many Analysts,
  Specification Curve, Forecasting Competition, Referee Report, Systematic Review (\S\ref{sec:research})\\
\bottomrule
\end{tabularx}
\end{center}

\section{Agents, skills and tasks}\label{sec:tasks}
Pollis knows three agents, Claude Code (\file{claude}), Codex (\file{codex}) and OpenCode (\code{opencode -{}-prompt}),
and any you add under \ui{Other Agents} on the Coding Agents page: a name and a command line, with \texttt{\small\{prompt\}}
where the first message goes. Without \texttt{\small\{prompt\}}, Pollis copies the first message to the clipboard, to paste
into the agent.

A \emph{task} is a menu item that starts an agent. Each has a \emph{tier}, the agent and model it runs on, and the
skills it reads. A tier is set for each task; the defaults are these.

\begin{center}\small
\renewcommand{\arraystretch}{1.2}
\begin{tabularx}{\linewidth}{@{}p{0.31\linewidth}p{0.17\linewidth}p{0.12\linewidth}L@{}}
\toprule
\textbf{Task} & \textbf{Menu} & \textbf{Tier} & \textbf{Skill}\\
\midrule
Polish with Agent & Reports & Light & polish-report\\
Write a Vignette & Replication & Standard & write-vignette\\
Prepare for Registration & Replication & Standard & prepare-method\\
Package a Toolbox & Compose & Standard & package-toolbox\\
Plan a Replication & Replication & Standard & plan-replication\\
Explain the Differences & Replication & Standard & explain-differences\\
Speed Up a Function & Profiling & Standard & speed-up-function\\
Tables from the Article & Replication & Strong & tables-from-article\\
Write a Method & Replication & Strong & write-method\\
Replicate an Article & Replication & Strong & replicate-article\\
Create a Panel & New Panels & Strong & create-panel\\
Create Panels from a Request & New Panels & Strong & create-panel\\
Translate Julia to Rust & Compilation Lab & Strong & translate-to-rust\\
\bottomrule
\end{tabularx}
\end{center}
A skill started from \menupath{Compose \gt Agent Skills} with \ui{Run with My Agent} is no task of the table, and runs
on the Standard tier.

\section{Tiers and cells}\label{sec:tiers}
Not every task needs the strongest model, and a task that reads numbers off a page or writes a method should not get
the cheapest. Three tiers come with Pollis:
\begin{itemize}
  \item \textbf{Light}: small tasks, on a fast, cheap or local model.
  \item \textbf{Standard}: everyday tasks.
  \item \textbf{Strong}: hard tasks, where a mistake costs most.
\end{itemize}

\subsection{The table of agents and models}
\menupath{Compose \gt Orchestration \gt Model Tasks and Routing} is a table. Its columns are the agents, and its rows
are the models: the cloud models, then the local models of this computer. A tier is a cell: type its name in a cell,
and the tier is that agent on that model. A tier is in one cell at most; a name of your own, such as \ui{Max}, makes
a new tier, which lasts while it is in a cell. A dash marks a model the agent cannot run. Until you change them, Light,
Standard and Strong are Claude Code on Claude Haiku, Sonnet and Opus.

\subsection{The cloud models}
Each cloud model is a row with, for each agent that runs it, the arguments that choose it. Pollis puts them after the
agent's program, so a tier's command line is the agent's, with the model chosen.

\begin{center}\small
\renewcommand{\arraystretch}{1.2}
\begin{tabularx}{\linewidth}{@{}p{0.22\linewidth}p{0.2\linewidth}p{0.17\linewidth}L@{}}
\toprule
\textbf{Model} & \textbf{Claude Code} & \textbf{Codex} & \textbf{OpenCode}\\
\midrule
Claude Haiku & \code{-{}-model haiku} & --- & \code{-{}-model anthropic/claude-haiku-5-5}\\
Claude Sonnet & \code{-{}-model sonnet} & --- & \code{-{}-model anthropic/claude-sonnet-5-5}\\
Claude Opus & \code{-{}-model opus} & --- & \code{-{}-model anthropic/claude-opus-5-5}\\
GPT-5.5 & --- & \code{-m gpt-5.5} & \code{-{}-model openai/gpt-5.5}\\
GPT-5.4 & --- & \code{-m gpt-5.4} & \code{-{}-model openai/gpt-5.4}\\
GPT-5.4 mini & --- & \code{-m gpt-5.4-mini} & \code{-{}-model openai/gpt-5.4-mini}\\
GPT-5.3 Codex & --- & \code{-m gpt-5.3-codex} & \code{-{}-model openai/gpt-5.3-codex}\\
Gemini 4 Argon & --- & --- & \code{-{}-model google/gemini-4-argon}\\
Gemini 3.8 Flash & --- & --- & \code{-{}-model google/gemini-3.8-flash}\\
Gemini 3.5 Flash-Lite & --- & --- & \code{-{}-model google/gemini-3.5-flash-lite}\\
\bottomrule
\end{tabularx}
\end{center}
The rows are the setting \setting{pollis.agents.cloudModels}: a model of your own is a row with its name and the
arguments of each agent, and an agent with none for a model cannot run it. Claude Code's \file{haiku}, \file{sonnet} and
\file{opus} are its names for the latest model of each kind.

\subsection{Choosing the tier of a task}
Below the table, \ui{Tasks} lists each task with its menu, its tier under \ui{Runs On}, and its skills. Choosing
another tier for a task overrides its default; a tier that leaves the table stops being an override, and the task goes
back to its default. Three settings keep it all:

\begin{center}\small
\renewcommand{\arraystretch}{1.2}
\begin{tabularx}{\linewidth}{@{}p{0.33\linewidth}L@{}}
\toprule
\textbf{Setting} & \textbf{What it holds}\\
\midrule
\file{pollis.agents.tierCells} & The cell of each tier: its \file{agent}, by name, and its \file{model},
  \file{cloud:} and a cloud model's name, or \file{local:} and a model file\\
\file{pollis.agents.cloudModels} & The cloud models of the rows, each with its arguments for each agent\\
\file{pollis.agents.taskTiers} & The tier of each task that is not on its default\\
\bottomrule
\end{tabularx}
\end{center}

\section{Local models}\label{sec:local}
A tier on a local model runs its agent on a model of this computer: free, private, and without a network. Pollis
serves the model with llama.cpp's \file{llama-server}, which comes with it, and points the agent at the server.
Claude Code, Codex and OpenCode can do this; an agent of your own cannot.

\begin{center}\small
\renewcommand{\arraystretch}{1.25}
\begin{tabularx}{\linewidth}{@{}p{0.17\linewidth}L@{}}
\toprule
\textbf{Agent} & \textbf{How it reaches the local model}\\
\midrule
Claude Code & Through the server's Anthropic API: \file{ANTHROPIC_BASE_URL} is the server, and the Opus, Sonnet and
  Haiku models are all the local model\\
OpenCode & Through the server's OpenAI API, as a provider \file{pollis-llamacpp} that \file{OPENCODE_CONFIG_CONTENT}
  describes\\
Codex & Through the server's OpenAI API, as the provider \file{pollis-llamacpp}, given with \code{-c} arguments\\
\bottomrule
\end{tabularx}
\end{center}
The agent's own choice of model, \code{-{}-model} or \code{-m}, is dropped, so the local model is the one used. The
model gets a context of 32,768 tokens, because an agent's instructions and tools alone take about 20,000. A server
already running for the model is used again.

\menupath{Compose \gt Local AI Models \gt Set Active Models} lists the local models: those on this computer, those
Pollis downloaded and those in the Hugging Face cache, with their kind, what each reads, whether it can call tools, and
the tiers that use it. An agent needs a chat model that can call tools.

\section{When a task starts}\label{sec:start}
\begin{enumerate}
  \item \textbf{The tier.} The task's tier, from \ui{Tasks}, and the cell of that tier, from the table. A tier in no
    cell stops here, with a notification that offers to open Model Tasks and Routing.
  \item \textbf{The agent.} The agent's program is looked for on this computer, as a terminal would find it. An agent
    that is not there stops here, with its install page.
  \item \textbf{The model.} A cloud model's arguments go after the program. A local model is started first, with a
    notification while llama.cpp loads it.
  \item \textbf{The terminal.} A terminal named \ui{Pollis Agents:} with the task and its tier, such as
    \ui{Pollis Agents: Write a Method (Strong)}, starts the agent in the task's folder, else in the workspace's first
    folder, else in your home folder.
  \item \textbf{The first message.} It points the agent at the skill and says what to do it on. It goes in the place of
    \texttt{\small\{prompt\}}, or into the clipboard.
\end{enumerate}
From then on you talk to the agent in its terminal. A coding agent reads and writes files and runs commands: read what
it proposes, and check its results, as you would a colleague's.

\section{Agents working together on research}\label{sec:research}
Some practices of empirical research need several people working on their own: many teams analysing the same data,
referees in different roles, two reviewers reading the same studies. One researcher alone usually cannot follow them;
with coding agents, one researcher can. Five pages of \menupath{Compose \gt Orchestration} are built on such
practices. They are in the menu, but not available yet.

\begin{center}\small
\renewcommand{\arraystretch}{1.25}
\begin{tabularx}{\linewidth}{@{}p{0.23\linewidth}LL@{}}
\toprule
\textbf{Page} & \textbf{What the agents do} & \textbf{The practice it follows}\\
\midrule
Many Analysts & Answer the same question on the same data, each on its own; Pollis shows how far their estimates
  spread & Many-analysts studies in psychology, economics and finance\\
Specification Curve & List the defensible choices in an analysis; Pollis runs every combination and shows the
  estimates in order & Specification curve analysis\\
Forecasting Competition & Each builds a forecasting model; the models are scored on data none of them saw, then
  combined & The M forecasting competitions, and the combination of forecasts\\
Referee Report & Referees in different roles read a draft paper, and an editor merges their reports & Peer review\\
Systematic Review & Two screen the studies and extract their results, each on its own; you settle where they
  differ, and Pollis runs the meta-analysis & Two independent reviewers\\
\bottomrule
\end{tabularx}
\end{center}

\subsection{Why several agents, and why different models}
Two analysts who make the same choices teach nothing about the choices. Agents on one model tend to choose alike, as
one person would on two days; agents on different models, from different providers, do not share all their habits.
Each research page lets you choose the agent and model of every participant, from the cells of Model Tasks and Routing,
and says so when the same one is chosen twice.

\subsection{Many Analysts}
Several agents get the same research question and the same dataset, and each analyses it on its own: its own
specification, its own controls, its own standard errors. Pollis gathers the estimates and shows how far they spread.

Different teams given the same data and the same question reach different answers, through the choices each makes.
Studies of many teams have shown it in psychology, \emph{Many Analysts, One Data Set} (2018); in economics, \emph{The
Influence of Hidden Researcher Decisions in Applied Microeconomics} (2021); and in finance, \emph{Nonstandard Errors}
(2024), with 164 teams, which found the spread across teams as large as the standard error each team reports, and
called it the \emph{nonstandard error}. A researcher cannot repeat such a study alone. The page uses what Pollis
already has: its datasets, the Julia packages, and the checking of results against the numbers reported.

The page will take four steps:
\begin{enumerate}
  \item \textbf{Question and Data.} The question in plain words; the data, a dataset of PollisDatasets or a file of your
    own; and the estimate every analyst must report, with its unit, such as \emph{the effect on 1978 earnings, in
    dollars}, so that the answers can be compared. The example is LaLonde's job training data: the experiment, and the
    comparison groups drawn from two surveys, whose estimates of the same effect spread widely.
  \item \textbf{Analysts.} From 3 to 12, each an agent on a model.
  \item \textbf{Runs.} All at once and unattended, each in a folder of its own, seeing only the question, the data and
    its folder. An analyst edits files without asking, inside its folder only.
  \item \textbf{Results.} Each estimate with its interval, in order; the median, the interquartile range and the
    nonstandard error, the standard deviation of the estimates across analysts, beside the median standard error the
    analysts report; and a table of each analyst's choices, which shows which choices move the estimate.
\end{enumerate}
Each analyst writes three files: \file{analysis.jl}, the whole analysis in Julia, from the data to the estimate;
\file{result.toml}, the estimate, its standard error, its interval and the number of observations, with the choices in
a few words each; and \file{REPORT.md}, the reasons for the choices. Pollis runs every \file{analysis.jl} again itself,
and marks an analysis whose result differs from the one it reports.

\subsection{Specification Curve}
Agents list the defensible choices in an analysis, such as which controls to include, which observations to drop and
which estimator to use. Pollis runs every combination of them and plots the estimates in order, with the choices of
each marked beneath it. This is the specification curve analysis of Simonsohn, Simmons and Nelson (2020). Where Many
Analysts shows the choices that independent analysts happen to make, the specification curve shows all of them, one
by one.

\subsection{Forecasting Competition}
Each agent builds a forecasting model for the same series. Pollis keeps the last part of the series from all of
them, scores every model on it, and then combines the forecasts. This follows the M competitions of forecasting
methods, and Bates and Granger (1969) on combining forecasts: a combination often forecasts better than each of the
models in it. In finance and economics, where a forecast is judged by the data that come after it, the page shows
which approach wins on the series at hand, and by how much.

\subsection{Referee Report}
Agents in different roles each read a draft paper:
\begin{itemize}
  \item an \textbf{econometrician}, on identification and inference;
  \item a \textbf{field expert}, on the literature;
  \item a \textbf{replicator}, who runs the paper's code again.
\end{itemize}
An \textbf{editor} agent then merges their notes into one report, with the points ranked by how much they matter.
It gives a report before submission, when changes are still cheap. The Knowledge Base and the references of Pollis
give the field expert sources to check the paper's claims against.

\subsection{Systematic Review}
This follows the standard of systematic reviews, two reviewers working independently:
\begin{itemize}
  \item two agents screen the abstracts against your criteria for including a study;
  \item two agents extract the effect sizes from the studies that pass;
  \item Pollis shows where they disagree, you decide, and Pollis then runs the meta-analysis.
\end{itemize}
The studies come from the references database and Literature Radar, which find and keep them.

\section{Under the surface}\label{sec:machinery}
\begin{itemize}
  \item \textbf{One module for the tiers.} \file{modelRouting.ts} holds the tasks, the tiers, the cloud models and the
    settings, and answers which command line a tier runs (\file{tierRun}). It imports no other module of Compose, so
    that every task can import it without a cycle: the agents are passed to it.
  \item \textbf{One function starts every task.} \file{startAgent} (\file{pollisAgents.ts}) takes the task's title,
    its first message, its folder and the task, finds the tier and its cell, looks for the program
    (\file{findProgram}), starts a local model when the cell has one, and opens the terminal with
    \file{waitOnExit}, so that what the agent last wrote stays on screen when it exits.
  \item \textbf{Local models.} \file{localModelLaunch} asks the Julia extension's command
    \file{language-julia.pollisLlamaServer} (\file{pollisLlamaServer.ts}) for a server of the model, and turns its
    address into the environment or the arguments of the agent.
  \item \textbf{The page.} \file{modelRoutingPage.ts} draws the table and the tasks, and writes the three settings.
  \item \textbf{The research pages.} They are greyed out in \file{compose.contribution.ts} until each is built. They
    will share one runner: unattended runs that end by themselves, an agent on a cell the page chooses, and one folder
    for each participant.
\end{itemize}

\section{Where things are}

\begin{center}\small
\renewcommand{\arraystretch}{1.2}
\begin{tabularx}{\linewidth}{@{}p{0.47\linewidth}L@{}}
\toprule
\textbf{File} & \textbf{Role}\\
\midrule
\file{.../compose/browser/modelRouting.ts} & The tasks, the tiers, the cloud models, the settings, \file{tierRun},
  \file{localModelLaunch}\\
\file{.../compose/browser/modelRoutingPage.ts} & The page Model Tasks and Routing\\
\file{.../compose/browser/pollisAgents.ts} & The agents, the skills, \file{startAgent}\\
\file{.../compose/browser/codingAgentPage.ts} & The page Coding Agents\\
\file{.../language-julia/src/pollisLlamaServer.ts} & llama.cpp's server for local models\\
\file{.../compose/browser/compose.contribution.ts} & The Orchestration menu\\
\bottomrule
\end{tabularx}
\end{center}

\end{document}
